\documentclass[10pt,a4paper]{amsart}
\usepackage[margin=19mm]{geometry}
\usepackage{amsmath,amssymb,mathtools}
\usepackage[scr=rsfs,cal=euler]{mathalfa}
\usepackage{xfrac}
\usepackage[unicode,hidelinks,bookmarks=false]{hyperref}
\newcommand{\ie}{i.e.\ }
\newcommand{\cf}{cf.\ }
\newcommand{\resp}{respectively\ }
\newcommand{\Subsection}{\subsection}
\providecommand{\nobreakdash}{\nobreak\hbox{-}}
\setlength{\emergencystretch}{2em}
\multlinegap0pt
\usepackage{algorithm}\usepackage{algpseudocode}
\newtheorem{lem}{Lemma}
\title{On a Boundary-Initial Value Problem for Fractional Differential Equation with Sequential Caputo derivatives}
\author{Fayziev Yusuf, Jumaeva Shakhnoza}
\date{Source: arXiv:2604.21319v1 (CC BY 4.0)}
\begin{document}
\maketitle

\noindent\emph{Source and licence.} On a Boundary-Initial Value Problem for Fractional Differential Equation with Sequential Caputo derivatives. Version arXiv:2604.21319v1 (CC BY 4.0), distributed by arXiv under the Creative Commons Attribution 4.0 International licence, http://creativecommons.org/licenses/by/4.0/, as recorded in the arXiv licence metadata for this exact version and shown on the abstract page https://arxiv.org/abs/2604.21319v1. Full licence notice: \texttt{licenses/arxiv-2604.21319-CC-BY-4.0.txt}.\par\smallskip

\noindent\textit{Editorial scope.} Counted unit: the article's lem jumaeva:eq:8 (source0.tex lines 309--350). The article's complete original proof of it is delivered with it (source0.tex lines 352--433, one \texttt{proof} environment of 82 lines). No mathematics is written, repaired or re-derived: the statement and the proof are the article's own lines, in the article's own notation, and no retained line is substituted or re-flowed.  Retained uncounted as the source-local context the proof needs: lines 267--269.  Nothing was omitted from any retained span, so the package carries no omission record.  No retained line contains a citation, so the wrapper carries no bibliography and no bibliography entry.  No retained line contains a figure environment, an image-inclusion command or drawing code, so no image is delivered and the package declares no asset dependency.  Editorial formatting only, in addition: an AMS \texttt{amsart} preamble that restates the article's own declarations and macros used by the retained lines, namely the environments \texttt{lem} on separate counters, together with the standard packages those lines need.  The retained environments print in this document as Lemma 1 in order of appearance, whereas the article prints the same items with the numbering of its own full document.  The complete native source of the article as distributed in the arXiv e-print for this exact version is bundled unchanged as \texttt{sources/199017/source0.tex}.
where the parameters and variables satisfy the conditions \begin{equation*}
\gamma_1, \gamma_2, \delta_1, \delta_2, \delta_3, x, y \in \mathbb{C} \quad \text{and} \quad \min\{\alpha_1, \alpha_2, \alpha_3, \alpha_4, \beta_1, \beta_2, \beta_3\} > 0.
\end{equation*} \label{jumaeva:def:3}

\begin{lem}
The following relations hold for $0<\alpha<1$:
\begin{enumerate}
    \item   \[
    \int_0^t \eta^{\delta_1-1 }E_2\left( \begin{array}{c}\gamma_1,\alpha_1, \beta_1 ; \gamma_2, \alpha_2 ; \\\delta_1, \alpha_3, \beta_2 ;\delta_2, \alpha_4; \delta_3, \beta_3 ; \end{array} \left| \begin{array}{c}\omega_1 \eta^{\alpha_3} \\\omega_2 \eta^{\beta_2} \end{array} \right.\right) d\eta
    \]
    \begin{equation} \label{jumaeva:eq:8}
    =t^{\delta_1} E_2\left( \begin{array}{c}\gamma_1,\alpha_1, \beta_1 ; \gamma_2, \alpha_2 ; \\\delta_1+1, \alpha_3, \beta_2 ;\delta_2, \alpha_4; \delta_3, \beta_3 ; \end{array} \left| \begin{array}{c}\omega_1 t^{\alpha_3} \\\omega_2 t^{\beta_2}\end{array}\right.\right).\end{equation}
    \item  
    \begin{equation*}
    D_t^\alpha\bigg[t^{\delta_1-1 }E_2\left(\begin{array}{c}\gamma_1,\alpha_1, \beta_1 ; \gamma_2, \alpha_2 ; \\\delta_1, \alpha_3, \beta_2 ;\delta_2, \alpha_4; \delta_3, \beta_3 ; \end{array} \left| \begin{array}{c}\omega_1 t^{\alpha_3} \\\omega_2 t^{\beta_2} \end{array} \right.\right)\bigg]
    \end{equation*}
    \begin{equation}\label{jumaeva:eq:9}
    =t^{\delta_1-\alpha-1} E_2\left( \begin{array}{c}\gamma_1,\alpha_1, \beta_1 ; \gamma_2, \alpha_2 ; \\\delta_1-\alpha, \alpha_3, \beta_2 ;\delta_2, \alpha_4; \delta_3, \beta_3 ; \end{array} \left|\begin{array}{c}\omega_1 t^{\alpha_3} \\\omega_2 t^{\beta_2} \end{array}\right.\right).
    \end{equation}
    \item  Let $\delta_1>1$. Then
    \begin{equation*}
    D_t^\alpha \bigg[\int_0^t \eta^{\delta_1-1} E_2\left(\begin{array}{c}
\gamma_1,\alpha_1, \beta_1 ; \gamma_2, \alpha_2 ; \\
\delta_1, \alpha_3, \beta_2 ;\delta_2, \alpha_4; \delta_3, \beta_3 ; 
\end{array} \left| 
\begin{array}{c}
\omega_1 \eta^{\alpha_3} \\
\omega_2 \eta^{\beta_2} 
\end{array} 
\right.
\right)f(t-\eta) d\eta \bigg] 
\end{equation*}
\begin{equation}\label{jumaeva:eq:10}
=\int_0^t \eta^{\delta_1-\alpha-1} E_2\left(\begin{array}{c}
\gamma_1,\alpha_1, \beta_1 ; \gamma_2, \alpha_2 ; \\
\delta_1-\alpha, \alpha_3, \beta_2 ;\delta_2, \alpha_4; \delta_3, \beta_3 ; 
\end{array} \left| 
\begin{array}{c}
\omega_1 \eta^{\alpha_3} \\
\omega_2 \eta^{\beta_2} 
\end{array} 
\right.
\right) f(t-\eta) d\eta.
\end{equation}
\end{enumerate}\label{jumaeva:lem:5} 
\end{lem} 

\begin{proof}
One may simply obtain the desired equalities \eqref{jumaeva:eq:8} and \eqref{jumaeva:eq:9}  using Definition \ref{jumaeva:def:3} and the definite integral of the power function. For proving \eqref{jumaeva:eq:10}, we initially use the definition of the Caputo fractional derivative:
\begin{align*}
   D_t^\alpha\big[F(t)\big]&=\frac{1}{\Gamma(1-\alpha)}\int_0^t \frac{F'(s)}{(t-s)^\alpha} ds= \\
   &=\frac{1}{\Gamma(1-\alpha)}\int_0^t \int_0^s \frac{f(\eta)}{(t-s)^\alpha}\frac{\partial}{\partial s}\bigg[(s-\eta)^{\delta_1-1 }E_2\left(\begin{array}{c}
\gamma_1,\alpha_1, \beta_1 ; \gamma_2, \alpha_2 ; \\
\delta_1, \alpha_3, \beta_2 ;\delta_2, \alpha_4; \delta_3, \beta_3 ; 
\end{array} \left| 
\begin{array}{c}
\omega_1 (s-\eta)^{\alpha_3} \\
\omega_2 (s-\eta)^{\beta_2} 
\end{array} 
\right.
\right)\bigg]d\eta ds,
\end{align*}
where $F(t)=\int_0^t (t-\eta)^{\delta_1-1} E_2\left(\begin{array}{c}
\gamma_1,\alpha_1, \beta_1 ; \gamma_2, \alpha_2 ; \\
\delta_1, \alpha_3, \beta_2 ;\delta_2, \alpha_4; \delta_3, \beta_3 ; 
\end{array} \left| 
\begin{array}{c}
\omega_1 (t-\eta)^{\alpha_3} \\
\omega_2 (t-\eta)^{\beta_2} 
\end{array} 
\right.
\right)f(\eta) d\eta. $

Using the term-wise differentiation rule for $E_2(\cdot,\cdot)$ and then swapping the order of integration yields for $\delta_1>1$:
\begin{equation*}
     D_t^\alpha\big[F(t)\big]=\frac{1}{\Gamma(1-\alpha)}\int_0^t \int_0^s \frac{f(\eta)}{(t-s)^\alpha}(s-\eta)^{\delta_1-2 }E_2\left(\begin{array}{c}
\gamma_1,\alpha_1, \beta_1 ; \gamma_2, \alpha_2 ; \\
\delta_1-1, \alpha_3, \beta_2 ;\delta_2, \alpha_4; \delta_3, \beta_3 ; 
\end{array} \left| 
\begin{array}{c}
\omega_1 (s-\eta)^{\alpha_3} \\
\omega_2 (s-\eta)^{\beta_2} 
\end{array} 
\right.
\right)d\eta ds
\end{equation*}
\begin{equation*}
 =\frac{1}{\Gamma(1-\alpha)}\int_0^t f(\eta) \int_\eta^t \frac{(s-\eta)^{\delta_1-2}}{(t-s)^\alpha} E_2\left(\begin{array}{c}
\gamma_1,\alpha_1, \beta_1 ; \gamma_2, \alpha_2 ; \\
\delta_1-1, \alpha_3, \beta_2 ;\delta_2, \alpha_4; \delta_3, \beta_3 ; 
\end{array} \left| 
\begin{array}{c}
\omega_1 (s-\eta)^{\alpha_3} \\
\omega_2 (s-\eta)^{\beta_2} 
\end{array} 
\right.
\right)ds d\eta. 
\end{equation*}
Now, we expand  the function $E_2(\cdot,\cdot)$ into its defining double series and integrate termwise. With
$x=\frac{s-\eta}{t-\eta}$, the inner integral becomes a Beta integral and produces the factor
$\Gamma(\alpha_3 m+\beta_2 k+\delta_1-1)/\Gamma(\alpha_3 m+\beta_2 k+\delta_1-\alpha)$:
\begin{equation*}
    D_t^\alpha\big[F(t)\big]=\frac{1}{\Gamma(1-\alpha)} \int_0^t f(\eta) \sum_{m=0}^\infty \sum_{k=0}^\infty \frac{\Gamma(\alpha_1 m+\beta_1 k+\gamma_1)\Gamma(\alpha_2 m+\gamma_2) \omega_1^m \omega_2^k}{\Gamma(\gamma_1)\Gamma(\gamma_2)\Gamma(\alpha_3 m+\beta_2 k+\delta_1-1)\Gamma(\alpha_4 m+\delta_2)\Gamma(\beta_3 k+ \delta_3)}
\end{equation*}
\begin{equation*}
    \times\int_0^1 (t-\eta)^{\delta_1-2} x^{\delta_1-2} (t-\eta)^{-\alpha}(1-x)^{-\alpha}(t-\eta)^{\alpha_3 m+\beta_2 k}x^{\alpha_3 m+\beta_2 k} (t-\eta) dx d\eta
\end{equation*}

\begin{equation*}
    =\frac{1}{\Gamma(1-\alpha)} \int_0^t f(\eta) \sum_{m=0}^\infty \sum_{k=0}^\infty \frac{\Gamma(\alpha_1 m+\beta_1 k+\gamma_1)\Gamma(\alpha_2 m+\gamma_2) (\omega_1(t-\eta)^{\alpha_3})^m (\omega_2(t-\eta)^{\beta_2})^k}{\Gamma(\gamma_1)\Gamma(\gamma_2)\Gamma(\alpha_3 m+\beta_2 k+\delta_1-\alpha)\Gamma(\alpha_4 m+\delta_2)\Gamma(\beta_3 k+ \delta_3)}
\end{equation*}
\begin{equation*}
\times    (t-\eta)^{\delta_1-\alpha-1} d\eta
\end{equation*}
Hence, we get 
\begin{equation*}
   D_t^\alpha\big[F(t)\big]=\int_0^t (t-\eta)^{\delta_1-\alpha-1} E_2\left(\begin{array}{c}
\gamma_1,\alpha_1, \beta_1 ; \gamma_2, \alpha_2 ; \\
\delta_1-\alpha, \alpha_3, \beta_2 ;\delta_2, \alpha_4; \delta_3, \beta_3 ; 
\end{array} \left| 
\begin{array}{c}
\omega_1 (t-\eta)^{\alpha_3} \\
\omega_2 (t-\eta)^{\beta_2} 
\end{array} 
\right.
\right) f(\eta) d\eta.
\end{equation*}
which is the claimed relation \eqref{jumaeva:eq:10}.
\end{proof}

\end{document}
